% ============================================================================
\section{Toward a Hybrid Architecture}
\label{sec:hybrid}
% ============================================================================

The comparison in the previous section suggests that no single paradigm
is optimal across all criteria. This motivates a \concept{hybrid
architecture} that combines the strengths of multiple paradigms. This
section sketches a working hypothesis for such an architecture,
targeting MNIST~\cite{lecun1998mnist} as a proof of concept.

\subsection{Design Principles}

Our hybrid design is guided by four principles:

\begin{enumerate}[leftmargin=2em]
  \item \textbf{Interpretable feature extraction.} Use Tsetlin
    Machine clauses for the first stage: they learn human-readable
    rules about which pixel patterns indicate which digits.
  \item \textbf{Flexible composition.} Use Logic Gate Networks for
    the middle layers: they can learn complex combinations of the
    clause outputs using arbitrary logic gates.
  \item \textbf{Efficient deployment.} Use LUT compilation for
    any continuous functions that remain: compile learned activations
    into tables so that inference is pure memory access + bitwise
    logic.
  \item \textbf{All in Rust.} Implement both training and inference
    in Rust for speed, safety, and cross-platform deployment.
\end{enumerate}

\subsection{Proposed Architecture}

\begin{figure}[htbp]
\centering
\begin{tikzpicture}[
  stage/.style={draw, rounded corners=5pt, minimum width=3.5cm,
    minimum height=2cm, align=center, font=\small\sffamily,
    line width=1pt},
  arr/.style={-{Stealth[length=6pt]}, thick},
  lbl/.style={font=\scriptsize\sffamily, bitgray}
]
  % Stage 1: Input
  \node[stage, fill=lightbg] (input) at (0, 0)
    {\textbf{Input}\\[3pt]28$\times$28 pixels\\$\downarrow$\\784 Boolean features\\+ 784 negations};

  % Stage 2: Tsetlin clauses
  \node[stage, fill=bitgreen!10, draw=bitgreen] (tm) at (5.5, 0)
    {\textbf{Tsetlin Clauses}\\[3pt]$K \times M$ clauses\\learn feature patterns\\output: clause votes};

  % Stage 3: LGN composition
  \node[stage, fill=bitorange!10, draw=bitorange] (lgn) at (11, 0)
    {\textbf{Logic Gate Network}\\[3pt]Learned gates combine\\clause outputs\\gate types: AND/OR/XOR/\ldots};

  % Stage 4: Output
  \node[stage, fill=bitpurple!10, draw=bitpurple] (out) at (5.5, -4)
    {\textbf{LUT Classification}\\[3pt]Compiled scoring\\functions\\output: digit 0--9};

  % Arrows
  \draw[arr] (input) -- (tm);
  \draw[arr] (tm) -- (lgn);
  \draw[arr] (lgn) |- (out);

  % Labels
  \node[lbl, above=0.1cm] at (2.75, 0.8)
    {booleanise};
  \node[lbl, above=0.1cm] at (8.25, 0.8)
    {clause $\to$ gate};
  \node[lbl, right=0.1cm] at (11.8, -2)
    {compile to LUTs};

  % Training annotations
  \draw[dashed, bitgreen, thick] (-1.8, -1.5) rectangle (7.5, 1.5);
  \node[font=\scriptsize\sffamily, bitgreen, anchor=south west]
    at (-1.8, 1.5) {Phase 1: Train with Tsetlin automata (no FP)};

  \draw[dashed, bitorange, thick] (4, -1.8) rectangle (13, 1.8);
  \node[font=\scriptsize\sffamily, bitorange, anchor=south east]
    at (13, 1.8) {Phase 2: Train gate types (backprop)};

  \draw[dashed, bitpurple, thick] (3.3, -5.5) rectangle (8, -2.5);
  \node[font=\scriptsize\sffamily, bitpurple, anchor=south west]
    at (3.3, -5.5) {Phase 3: Compile to LUTs};
\end{tikzpicture}
\caption{Proposed hybrid architecture. Stage 1 (Tsetlin) learns
  interpretable feature patterns. Stage 2 (LGN) composes patterns with
  learned logic gates. Stage 3 (LUT) compiles any remaining scoring
  functions into lookup tables. Training is phased: Tsetlin automata
  first (no floating-point needed), then backpropagation for gate
  learning, then one-time LUT compilation.}
\label{fig:hybrid-architecture}
\end{figure}

\subsection{Training Pipeline}

The hybrid architecture suggests a three-phase training pipeline:

\begin{description}[leftmargin=1em]
  \item[Phase 1: Clause learning (Tsetlin).]
    Train the Tsetlin Machine component on the booleanised MNIST
    inputs. This produces a set of interpretable clauses---Boolean
    expressions that capture local pixel patterns. The output of this
    phase is a fixed set of clause activations per input image.
    \emph{No floating-point arithmetic required.}

  \item[Phase 2: Gate learning (LGN).]
    Freeze the Tsetlin clauses. Train a Logic Gate Network on top
    of the clause outputs. The LGN learns which gate types and
    connectivity best combine the clause features into class scores.
    This phase uses standard backpropagation with softmax relaxation
    and temperature annealing.

  \item[Phase 3: LUT compilation.]
    After gate types are committed, compile any remaining continuous
    functions (e.g., scoring aggregation, temperature scaling) into
    lookup tables. The final inference model is entirely
    integer/Boolean.
\end{description}

\begin{engineernote}
This phased approach has a practical advantage for Rust
implementation: Phase~1 needs only integer operations and random
number generation. Phase~2 needs floating-point backpropagation
(the most complex part). Phase~3 is a one-time offline step.
You can implement and test each phase independently.
\end{engineernote}

\subsection{Expected Characteristics}

\begin{table}[htbp]
\centering
\caption{Expected characteristics of the hybrid architecture on MNIST
  (estimates based on component performance).}
\label{tab:hybrid-expected}
\begin{tabular}{@{}lr@{}}
  \toprule
  \textbf{Metric} & \textbf{Expected Range} \\
  \midrule
  Accuracy          & 97--99\% \\
  Model size         & 50--500\,kB \\
  Inference ops      & Boolean + LUT only \\
  FP at inference?   & No \\
  Interpretability   & Partial (Tsetlin clauses readable) \\
  \bottomrule
\end{tabular}
\end{table}

\subsection{Benchmarking Plan}

To validate the hybrid approach, we will benchmark against a standard
floating-point baseline:

\begin{enumerate}[leftmargin=2em]
  \item \textbf{Baseline model:} A standard 2-layer MLP with ReLU
    activations, trained with Adam optimiser in 32-bit floating
    point. Expected accuracy: $\sim$98\%.
  \item \textbf{Metrics:}
    \begin{itemize}[leftmargin=1em]
      \item Classification accuracy (test set)
      \item Inference latency (microseconds per image)
      \item Model size (bytes)
      \item Inference energy estimate (based on operation counts)
    \end{itemize}
  \item \textbf{Environment:} Both models implemented in Rust, running
    on the same CPU, with no GPU acceleration for either.
\end{enumerate}
